% !PRINTED 595
\weberpage{595}
%<<CONT
and from (5) there results, by squaring and making use of (4):
\[
\pi^2 \frac{4\psi(\boldsymbol{s})^4 - \lambda\varphi(\boldsymbol{s})^2\psi(\boldsymbol{s})^2 + 4\psi(\boldsymbol{s})^4}{4 - \lambda \boldsymbol{s}^2 + 4\boldsymbol{s}^4} = [\varphi'(\boldsymbol{s})\psi(\boldsymbol{s}) - \varphi'(\boldsymbol{s})\psi(\boldsymbol{s})]^2. % [sic] German source prints both the repeated numerator term and repeated right-hand factor; (5) and the next display give the intended complementary terms
\]
Since on the right there stands an integral function of $\boldsymbol{s}$, the numerator on the left-hand side must be divisible by the denominator (which also follows from consideration of the zero points), and it follows:
\[
\pi^2 \Phi = [\varphi'(\boldsymbol{s})\psi(\boldsymbol{s}) - \varphi(\boldsymbol{s})\psi'(\boldsymbol{s})]^2,
\]
where $\Phi$ is an integral function of $\boldsymbol{s}$ whose coefficients are integers in $\mathfrak{L}$. Hence $\Phi$ must be the square of an integral function whose coefficients are likewise integral algebraic numbers (Bd. I, § 2; Bd. II, § 159), and from this there results, for a variable $x$:
\weq{6}{\varphi'(x)\psi(x) - \psi'(x)\varphi(x) \equiv 0 \ (\mathrm{mod}\ \pi).}
From this it can be proved that the coefficients $A_1, A_3, A_5, \ldots, A_{p-2}$ are all divisible by $\pi$.

Of $A_1$ we know this already. Let us then assume that it has been proved for
\[
A_1, A_3, \ldots, A_{2\nu-1} \quad (2\nu+1 < p),
\]
then it follows:
\[
\begin{aligned}
\varphi(x) &\equiv A_{2\nu+1}x^{2\nu+1} + \cdots\\
\varphi'(x) &\equiv (2\nu+1)A_{2\nu+1}x^{2\nu} + \cdots \qquad (\mathrm{mod}\ \pi)\\
\psi(x) &\equiv 1 + A_{p-2}x^2 + \cdots\\
\psi'(x) &\equiv 2A_{p-2}x + \cdots
\end{aligned}
\]
If we now seek in (6) the coefficient of $x^{2\nu}$, it follows:
\[
(2\nu+1)A_{2\nu+1} \equiv 0 \ (\mathrm{mod}\ \pi).
\]
Since $2\nu+1 < p$ and $p$ is a prime, $2\nu+1$ is relatively prime to $\pi$, and it follows:
\weq{7}{A_{2\nu+1} \equiv 0 \ (\mathrm{mod}\ \pi).}
If in (4) we put for $u$ one of the values $2\alpha\xi/\mu$, § 158, (9), then $\boldsymbol{s}(u)$ becomes a root of the ideal division equation § 158, (10), and it follows:
\[
\pm \boldsymbol{s}(\pi u) \equiv [\boldsymbol{s}(u)]^p \ (\mathrm{mod}\ \pi),
\]
and if we raise to a square:
\weq{8}{\boldsymbol{s}(\pi u)^2 \equiv [\boldsymbol{s}(u)]^{2p} \ (\mathrm{mod}\ \pi).}
In particular, therefore, we have
\weq{9}{\boldsymbol{s}(u)^2 \equiv \boldsymbol{s}(u)^{2p} \ (\mathrm{mod}\ \pi),}
if
\weq{10}{\pi \equiv \pm 1 \ (\mathrm{mod}\ \mathfrak{m})}
%CONT>>
